\documentclass[a3paper, landscape, 11pt, twocolumn]{article}

% 引入宏包
\usepackage{fontspec}
\usepackage[UTF8]{ctex}
\usepackage{geometry}
\usepackage{listings}
\usepackage{xcolor}
\usepackage{enumitem}
\usepackage{tabularx}
\usepackage{makecell}
\usepackage{fancyhdr}
\usepackage{amssymb}
\usepackage{lastpage}
\usepackage{titlesec}
\usepackage{multicol}

% 设置页面边距 (A3横向双栏)
\geometry{left=2cm, right=2cm, top=2.5cm, bottom=2.5cm, columnsep=1.8cm}

% 颜色定义
\definecolor{codebg}{RGB}{245,245,245}
\definecolor{keywordcolor}{RGB}{0,0,150}
\definecolor{commentcolor}{RGB}{0,128,0}
\definecolor{stringcolor}{RGB}{163,21,21}
\definecolor{answerred}{RGB}{255,0,0} % 答案红色

% 页眉页脚设置
\pagestyle{fancy}
\fancyhf{}
\fancyhead[C]{\Large\bfseries 《面向对象程序设计》简答题专项练习（含详细答案）}
\fancyfoot[C]{第 \thepage\ 页 共 \pageref{LastPage}\ 页}
\fancyfoot[R]{得分：\underline{\hspace{2cm}}}
\fancyfoot[L]{题号：\underline{\hspace{1cm}}}
\renewcommand{\headrulewidth}{0.4pt}

% 简答题序号环境
\newcounter{shortq}
\setcounter{shortq}{0}
\newcommand{\sq}{%
    \refstepcounter{shortq}%
    \noindent\textbf{\theshortq、}\hspace{0.5em}%
}

% 答案命令（红色，增加段落间距）
\newcommand{\answer}[1]{%
    \par\noindent{\color{answerred} 答：#1}\par\vspace{0.8cm}%
}


% 继续答案命令（无“答：”，用于代码块打断后）
\newcommand{\answercontinue}[1]{%
    \par\noindent{\color{answerred} #1}\par\vspace{0.8cm}%
}

\begin{document}

% 试卷头部
\twocolumn[{
 \begin{center}
 \vspace*{-1cm}
 {\huge\bfseries 《面向对象程序设计》简答题专项练习（含详细答案）}\\[0.8em]
%  {\large 考试时间：120 分钟 \hspace{2em} 总分：150 分}\\[1em]
 {\large 姓名：\underline{\hspace{4cm}} \hspace{2em} 学号：\underline{\hspace{5cm}} \hspace{2em} 班级：\underline{\hspace{4cm}}}
 \vspace{0.8em}
 \hrule
 \vspace{1em}
 \end{center}
}]

\section*{【简答题】共 30 小题，每小题 5 分，共 150 分。}
\textit{（请将答案写在题目下方的空白处，要求要点明确、条理清晰）}

% --- 题目及参考答案（详细版）---

\sq 简述 Java 语言的主要特点。
\answer{Java 语言的主要特点包括：①平台无关性——Java 程序被编译为字节码，可在任何安装了 Java 虚拟机（JVM）的平台上运行，实现"一次编写，到处运行"；②面向对象——Java 支持封装、继承、多态等面向对象特性，使程序设计更加灵活和可维护；③自动内存管理——Java 提供垃圾回收机制（Garbage Collection），自动回收不再使用的对象，减轻程序员负担；④多线程支持——Java 内置多线程机制，支持并发程序设计；⑤安全稳定——Java 去除了指针操作，提供异常处理机制，保障程序运行的稳定性；⑥丰富的类库——Java 提供了大量标准类库（如集合框架、输入输出、网络编程等），方便快速开发。}

\sq 解释 JDK、JRE 和 JVM 的区别与联系。
\answer{JDK（Java Development Kit）是 Java 开发工具包，包含 JRE 以及 Java 开发工具（如编译器 javac、调试器、文档生成器等），用于 Java 程序的开发、编译和调试。JRE（Java Runtime Environment）是 Java 运行时环境，包含 JVM 和 Java 核心类库，用于运行已编译的 Java 程序。JVM（Java Virtual Machine）是 Java 虚拟机，负责执行 Java 字节码，是 Java 实现平台无关性的核心。三者的关系是：JDK 包含 JRE，JRE 包含 JVM。即开发 Java 程序需要 JDK，运行 Java 程序只需要 JRE，而 JVM 是 JRE 的核心组成部分。}

\sq 什么是标识符？Java 对标识符的命名规则有哪些？
\answer{标识符是用来标识变量、方法、类、接口等程序元素的名称。Java 标识符的命名规则：①必须以字母（包括中文、英文等 Unicode 字符）、下划线（\_）或美元符号（\$）开头，不能以数字开头；②后续字符可以是字母、数字、下划线或美元符号；③不能使用 Java 关键字或保留字（如 public、class、void 等）；④标识符区分大小写，例如 \texttt{name}和 \texttt{Name}是不同的标识符；⑤理论上没有长度限制，但建议见名知意，采用驼峰命名法（如 studentName、getAge）。}

\sq Java 的基本数据类型有哪些？并说明每种类型的占用字节数。
\answer{Java 共有 8 种基本数据类型：①byte（字节型），占用 1 个字节，取值范围 -128~127；②short（短整型），占用 2 个字节，取值范围 -32768~32767；③int（整型），占用 4 个字节，约±21 亿；④long（长整型），占用 8 个字节，需在数值后加 L 或 l；⑤float（单精度浮点型），占用 4 个字节，需在数值后加 F 或 f；⑥double（双精度浮点型），占用 8 个字节，默认浮点类型；⑦char（字符型），占用 2 个字节，采用 Unicode 编码，可存储任何字符；⑧boolean（布尔型），大小未明确定义（通常视为 1 字节），取值只有 true 和 false。}

\sq 简述自动类型转换和强制类型转换的区别，并举例说明。
\answer{自动类型转换（隐式转换）是指当将取值范围小的类型赋值给取值范围大的类型时，Java 自动完成转换，无需显式编写代码。转换方向为：byte→short→int→long→float→double，以及 char→int。例如：int i = 10; double d = i; // 自动将 int 转换为 double。强制类型转换（显式转换）是指将取值范围大的类型赋值给取值范围小的类型时，必须手动进行转换，可能会丢失精度。语法为：(目标类型)表达式。例如：double d = 3.14; int i = (int)d; // 结果为 3，小数部分丢失。强制转换还可能溢出，如将超出 byte 范围的值强制转为 byte，会得到截断后的结果。}

\sq 解释 \texttt{==} 和 \texttt{equals()} 方法的区别。
\answer{\texttt{==}是运算符，用于比较两个操作数的值是否相等。对于基本数据类型，\texttt{==}比较的是数值；对于引用类型，\texttt{==}比较的是内存地址，即判断两个引用是否指向同一个对象。\texttt{equals()}是 Object 类中定义的方法，用于比较两个对象的内容是否相等。默认实现与 \texttt{==}相同，比较地址，但许多类（如 String、包装类、集合类）重写了 \texttt{equals()}方法，使其比较对象的内容。例如：String s1 = new String("abc"); String s2 = new String("abc"); s1 == s2 返回 false（地址不同），s1.equals(s2)返回 true（内容相同）。因此，比较字符串内容应使用 equals()方法。}

\sq 数组的 \texttt{length} 属性和字符串的 \texttt{length()} 方法有何不同？
\answer{数组的 \texttt{length}是一个属性（成员变量），直接获取数组的长度（即数组中元素的个数）。例如：int[] arr = new int[5]; System.out.println(arr.length); // 输出 5。字符串的 \texttt{length()}是一个方法，返回字符串中字符的个数。例如：String str = "Hello"; System.out.println(str.length()); // 输出 5。这种差异源于 Java 的设计：数组是特殊对象，其长度通过公有属性 \texttt{length}访问；而 String 是普通类，长度通过方法 \texttt{length()}获取。需要注意的是，对于二维数组，\texttt{array.length}获取的是行数，每行的长度可通过 \texttt{array[i].length}获取。}

\sq 什么是类？什么是对象？二者之间的关系是什么？
\answer{类（Class）是对具有相同属性和行为的事物的抽象描述，是创建对象的模板。类定义了对象应该包含的成员变量（属性）和成员方法（行为）。例如，可以定义一个"学生"类，包含学号、姓名等属性和学习、考试等方法。对象（Object）是类的具体实例，是真实存在的实体。每个对象都拥有类中定义的属性和方法，但每个对象的属性值可以不同。例如，根据"学生"类可以创建出"张三"和"李四"两个具体的学生对象。类与对象的关系：类是对象的抽象（概念上的定义），对象是类的具体化（内存中的实体）。可以从一个类创建多个对象，如同从一个模具中生产出多个产品。}

\sq 简述构造方法的特点和作用。构造方法可以重载吗？
\answer{构造方法（Constructor）是类中一种特殊的方法，具有以下特点：①方法名必须与类名完全相同；②没有返回值类型，也不能用 void 声明；③在创建对象时（使用 new 关键字）自动调用；④每个类至少有一个构造方法，如果未定义，系统会提供一个无参的默认构造方法；⑤构造方法不能被继承，也不能被重写，但可以重载。构造方法的主要作用是：创建对象并为对象的成员变量进行初始化，使对象创建时就具有合理的初始状态。构造方法可以重载（Overload），即一个类中可以定义多个参数列表不同的构造方法，以便以不同方式初始化对象。例如，可以定义无参构造方法、带一个参数的构造方法和带多个参数的构造方法，分别实现不同的初始化逻辑。}

\sq \texttt{this} 关键字有哪些用途？请举例说明。
\answer{\texttt{this}关键字在 Java 中有以下几种主要用途：①调用本类中的成员变量，特别是当成员变量与局部变量同名时，用 \texttt{this.变量名}明确表示访问的是成员变量。例如：\texttt{this.name = name;}（将参数 name 赋值给成员变量 name）。②调用本类中的其他构造方法，语法为 \texttt{this(参数列表)}，且必须放在构造方法的第一行。这可以避免代码重复，实现构造方法间的调用。例如：一个构造方法中调用 \texttt{this(name, 0);}。③返回当前对象的引用，例如在方法中 \texttt{return this;}可以用于实现链式调用（如 \texttt{obj.setName("Tom").setAge(18);}）。④将当前对象作为参数传递给其他方法或构造器。\texttt{this}只能在非静态方法中使用，因为静态方法中不存在当前对象的概念。}

\sq 什么是方法重载？方法重载的规则是什么？
\answer{方法重载（Overload）是指在同一类中，允许存在多个同名的方法，只要它们的参数列表不同即可。方法重载提供了以统一方式处理不同类型或不同数量参数的能力，增强了程序的灵活性和可读性。例如，\texttt{System.out.println()}方法就有多个重载版本，可以接受不同类型的数据。方法重载必须遵循的规则：①方法名必须相同；②参数列表必须不同（参数个数不同、参数类型不同、参数顺序不同，三者满足其一即可）；③返回值类型可以相同也可以不同，但不能仅以返回值类型不同作为重载的依据；④访问修饰符可以任意，不影响重载；⑤可以抛出不同的异常。重载方法在编译时根据传入参数的类型和数量确定调用哪个版本，属于编译时多态。}

\sq 什么是方法重写？方法重写需要遵循哪些规则？
\answer{方法重写（Override）是指子类根据自身需要，重新定义从父类继承来的某个方法。重写后，子类对象调用该方法时将执行子类中定义的新版本，而不是父类的版本。方法重写是实现运行时多态的基础。方法重写必须遵循的规则：①方法名、参数列表必须与被重写的方法完全相同；②返回值类型可以是原返回值类型的子类型（协变返回类型）；③访问权限不能比父类中被重写的方法更严格（例如父类方法是 public，子类不能改为 private）；④不能抛出比父类方法更宽泛的异常（即子类方法抛出的异常必须是父类方法抛出异常的子类或不抛出异常）；⑤被重写的方法不能是 final 类型的（final 方法禁止重写）；⑥被重写的方法不能是 static（静态方法属于类，不存在重写概念，只能隐藏）；⑦构造方法不能被重写。}

\sq 简述类变量（静态变量）和实例变量的区别。
\answer{类变量（静态变量）是指使用 \texttt{static}关键字修饰的成员变量，它属于整个类，而不属于某个特定的对象。实例变量是没有 \texttt{static}修饰的成员变量，每个对象都有自己独立的一份副本。具体区别如下：①存储位置：类变量存储在方法区（或堆中，取决于 JDK 版本），实例变量存储在堆内存中每个对象自己的空间内。②生命周期：类变量在类加载时分配内存，在类卸载时释放；实例变量在创建对象时分配，在对象被垃圾回收时释放。③访问方式：类变量可以通过类名直接访问（推荐），也可以通过对象访问；实例变量只能通过对象访问。④共享性：类变量被该类的所有对象共享，修改类变量会影响所有对象；实例变量为每个对象私有，互不影响。⑤所属关系：类变量属于类，实例变量属于对象。}

\sq 静态方法和实例方法有何不同？静态方法能否访问实例变量？为什么？
\answer{静态方法（使用 \texttt{static}修饰的方法）和实例方法的区别：①归属不同：静态方法属于类，可以通过类名直接调用；实例方法属于对象，必须通过对象引用调用。②访问权限：静态方法只能直接访问静态成员（静态变量和其他静态方法），不能直接访问实例变量和实例方法；实例方法可以直接访问静态成员和实例成员。③\texttt{this}关键字：静态方法中没有 \texttt{this}引用（因为不依赖于具体对象），实例方法中有 \texttt{this}引用指向当前对象。④内存分配：静态方法在类加载时分配入口地址；实例方法在创建对象时分配入口地址（但方法代码本身在方法区共享，只是调用时传入不同对象）。静态方法不能直接访问实例变量的原因：静态方法在类加载时就已存在，而此时可能还没有任何对象被创建，因此无法确定要访问哪个对象的实例变量。如果需要在静态方法中访问实例变量，必须通过创建对象或传入对象引用的方式间接访问。}

\sq \texttt{final} 关键字可以修饰哪些地方？分别起什么作用？
\answer{\texttt{final}关键字在 Java 中可用于修饰类、方法和变量，各自作用如下：①修饰类——该类被称为最终类，表示该类不能被继承。例如，\texttt{public final class String { ... }}，String 类被设计为 final，防止被继承修改其不可变性。②修饰方法——该方法被称为最终方法，表示该方法不能被子类重写（Override）。这通常用于保护那些不希望子类改变的关键方法。注意：final 方法可以被继承，但不能被修改。③修饰变量——该变量成为常量，一旦赋值后不能再被改变。对于基本类型变量，值不可变；对于引用类型变量，引用不可变（即不能再指向其他对象），但对象内部状态可以改变。final 成员变量必须在声明时、构造方法中或构造代码块中完成初始化。final 局部变量可以在使用时初始化一次。}

\sq 什么是访问控制修饰符？Java 中的访问控制修饰符有哪些？各自的作用范围是什么？
\answer{访问控制修饰符用于控制类、成员变量、方法等程序元素的可见性和访问权限，是 Java 封装特性的重要体现。Java 提供了四种访问控制级别，按权限从小到大排列：①private（私有）——访问权限最小，只能在同一个类内部访问，对外完全隐藏。用于封装类的内部实现细节。②default（默认，无修饰符）——包级私有，只能被同一个包中的类访问。如果在定义成员时未加任何访问修饰符，则采用此级别。③protected（受保护）——允许被同一个包中的类访问，也允许被不同包中的子类访问（通过继承关系）。④public（公有）——访问权限最大，可以被任何类访问（只要类本身可访问）。访问控制修饰符有助于实现信息隐藏和模块化设计，提高代码安全性和可维护性。}

\sq 简述封装的概念及其在 Java 中的实现方式。
\answer{封装（Encapsulation）是面向对象三大特性之一，指将对象的属性（数据）和行为（方法）捆绑在一起，并对外部隐藏内部实现细节，仅对外提供公共的访问接口。封装的核心思想是"信息隐藏"，即对象的状态（成员变量）只能通过对象自己的方法来改变和读取，外部不能直接修改对象的内部状态。在 Java 中，封装主要通过以下方式实现：①使用访问控制修饰符（通常用 private）隐藏成员变量，防止外部直接访问；②为每个需要被外部访问的私有变量提供 public 的 getter 和 setter 方法，通过这些方法来读取和修改变量值；③在 getter/setter 中可以添加数据验证逻辑，确保数据的有效性（如年龄不能为负数）；④类内部的其他方法也可以根据需要封装复杂逻辑，仅对外提供简单接口。封装的好处包括：提高代码安全性、降低耦合度、便于维护和扩展。}

\sq 什么是继承？Java 中如何实现继承？子类能从父类继承哪些成员？
\answer{继承（Inheritance）是面向对象三大特性之一，指子类自动获得父类的属性和方法，并可以在其基础上扩展新的属性和方法。继承体现了"is-a"关系，实现了代码重用，提高了软件的可维护性和可扩展性。在 Java 中，使用 \texttt{extends}关键字实现继承，语法为：\texttt{class 子类名 extends 父类名 { ... }}。Java 只支持单继承，即一个子类只能有一个直接父类，但可以有多层继承（如 A→B→C）。子类能从父类继承的成员包括：①public 和 protected 修饰的成员（无论是否同包）；②同包中默认访问权限的成员；③父类中所有非 private 成员变量和方法（但私有成员不能被直接访问，只能通过父类提供的公有方法间接访问）。子类不能继承父类的构造方法（但可以通过 \texttt{super()}调用），也不能继承父类的私有成员和不同包中的默认访问权限成员。}

\sq \texttt{super} 关键字的作用是什么？请举例说明。
\answer{\texttt{super}关键字用于在子类中引用父类的成员，主要有以下几种用途：①调用父类的构造方法——在子类构造方法的第一行使用 \texttt{super(参数列表)}，可以调用父类对应的构造方法完成父类部分的初始化。如果子类构造方法中没有显式调用 \texttt{super()}，编译器会自动插入无参的 \texttt{super()}（前提是父类有无参构造方法）。②访问父类被隐藏的成员变量——当子类定义了与父类同名的成员变量时，父类的成员变量被隐藏，可以通过 \texttt{super.变量名}访问父类的变量。③调用父类被重写的方法——当子类重写了父类的方法，但仍想调用父类原来的实现时，可以使用 \texttt{super.方法名()}。例如：\texttt{super.display();}。注意：\texttt{super}不能用于静态上下文（静态方法中）。\texttt{super}和 \texttt{this}有相似之处，但 \texttt{this}指向当前对象，\texttt{super}指向父类对象部分。}

\sq 什么是多态？Java 中多态的实现方式有哪些？
\answer{多态（Polymorphism）是面向对象三大特性之一，指同一种操作作用于不同的对象，可以产生不同的执行结果。简单说，就是"一种接口，多种实现"。多态提高了程序的灵活性和可扩展性。Java 中多态的实现方式分为两类：①编译时多态——通过方法重载（Overload）实现。在编译阶段，根据方法调用的参数个数、类型和顺序确定要执行的具体方法。②运行时多态——通过方法重写（Override）和向上转型（父类引用指向子类对象）实现。当父类引用调用被子类重写的方法时，JVM 会根据引用实际指向的对象类型决定执行哪个版本的方法，这个过程称为动态绑定（Dynamic Binding）。运行时多态需要满足三个条件：继承（或实现接口）、方法重写、父类引用指向子类对象。多态的优点包括：提高代码的复用性、降低耦合度、便于扩展和维护。}

\sq 什么是抽象类？抽象类有哪些特点？
\answer{抽象类是用 \texttt{abstract}关键字修饰的类，它是一种不能被实例化（不能创建对象）的类，通常作为其他类的基类（父类），用于定义一些通用的属性和行为规范。抽象类的特点如下：①抽象类必须用 \texttt{abstract}修饰；②抽象类不能实例化，只能被继承；③抽象类可以包含抽象方法和非抽象方法；④抽象类可以有构造方法，用于子类实例化时初始化父类成员；⑤抽象类可以包含成员变量、常量、静态成员等；⑥如果子类继承抽象类，则子类必须实现父类中所有的抽象方法，除非子类也是抽象类；⑦抽象类中可以不包含任何抽象方法，但包含抽象方法的类必须是抽象类；⑧抽象类不能被 \texttt{final}修饰（因为 final 类不能继承，与抽象类设计初衷矛盾）。抽象类主要用于定义标准、规范，强制子类实现特定的方法。}

\sq 什么是接口？接口和抽象类的区别是什么？
\answer{接口（Interface）用 \texttt{interface}关键字定义，是一种完全抽象的类（JDK8 之前），只包含常量和抽象方法的定义（JDK8 开始允许默认方法和静态方法，JDK9 开始允许私有方法）。接口用于定义一组规范，指明类应该提供哪些功能，但不关心如何实现。接口和抽象类的区别如下表所示：

| 比较维度 | 抽象类 | 接口 |
|---------|-------|------|
| 关键字 | abstract class | interface |
| 继承/实现 | 单继承 | 多实现（类可实现多个接口） |
| 构造方法 | 可以有构造方法 | 不能有构造方法 |
| 成员变量 | 可以是各种变量 | 默认 public static final 常量 |
| 方法类型 | 抽象、非抽象方法均可 | 抽象方法、默认方法、静态方法、私有方法 |
| 访问修饰符 | 可自由指定 | 默认 public |
| 设计思想 | is-a 关系（模板） | like-a 关系（能力/规范） |
| 添加新方法 | 不影响子类 | 需提供默认实现，否则影响实现类 |

两者联系：接口可以继承多个接口，抽象类可以实现接口。}

\sq Java 中如何实现多重继承？为什么？
\answer{Java 不支持类的多重继承（即一个类不能直接继承多个父类），原因是多重继承可能带来"菱形问题"（Diamond Problem）——当两个父类包含相同的方法时，子类无法确定继承哪个版本，导致二义性和复杂性。例如，C++ 支持多重继承，就需要通过虚继承等复杂机制解决冲突。但 Java 提供了替代方案：通过接口实现多重继承的效果。一个类可以同时实现多个接口（用 \texttt{implements}关键字，接口间用逗号分隔），从而获得多个接口中定义的行为规范。接口中的方法都是抽象的，不存在具体实现冲突；如果多个接口有默认方法（default method）冲突，子类必须重写该方法解决冲突。这种方式既保留了多重继承的灵活性（一个类具备多种能力），又避免了复杂性。因此，Java 设计者认为"单继承 + 多实现"是更安全、更简洁的折中方案。}

\sq 什么是异常？Java 的异常处理机制中有哪些关键字？分别说明其作用。
\answer{异常（Exception）是程序运行过程中出现的非正常事件，它会中断程序的正常执行流程。例如，除数为 0、数组下标越界、文件不存在等都会引发异常。Java 提供了强大的异常处理机制，使用以下五个关键字：①try——用于监控可能抛出异常的代码块，try 块后面必须跟至少一个 catch 或 finally 块。②catch——用于捕获 try 块中抛出的异常，并根据异常类型进行相应处理。可以有多个 catch 块，按从上到下顺序匹配。③finally——无论是否发生异常，finally 块中的代码都会执行（除非 JVM 退出）。通常用于释放资源（如关闭文件、数据库连接）。④throw——用于手动抛出一个异常对象（可以是自定义异常），一般用于方法内部。⑤throws——用于在方法声明中说明该方法可能抛出的异常类型，告知调用者需要处理这些异常。Java 异常处理机制将正常业务逻辑和错误处理代码分离，提高了程序的可读性和健壮性。}

\sq 简述受检查异常（checked exception）和运行时异常（unchecked exception）的区别。
\answer{Java 中的异常分为两大类：受检查异常（Checked Exception）和运行时异常（Runtime Exception），主要区别如下：①定义与继承关系：受检查异常继承自 \texttt{Exception}（但不包括 \texttt{RuntimeException}及其子类）；运行时异常继承自 \texttt{RuntimeException}。②编译时检查：受检查异常必须在代码中显式处理（使用 try-catch 捕获或在方法上声明 throws），否则编译不通过；运行时异常不强制处理，编译器允许通过，通常是由程序逻辑错误引起（如空指针、除零、下标越界）。③典型例子：受检查异常如 \texttt{IOException}、\texttt{SQLException}、\texttt{ClassNotFoundException}；运行时异常如 \texttt{NullPointerException}、\texttt{ArithmeticException}、\texttt{ArrayIndexOutOfBoundsException}。④处理建议：受检查异常表示外部环境问题，调用者应该处理；运行时异常表示程序内部错误，应通过改进代码避免。Java 设计者认为，对于那些可预见的、调用者可以恢复的异常，应设计为受检查异常。}

\sq 自定义异常类应如何定义？请写出基本步骤。
\answer{自定义异常类用于描述特定业务场景下的错误情况，使异常处理更具语义性。定义自定义异常的基本步骤如下：①选择父类——通常继承 \texttt{Exception}（成为受检查异常）或继承 \texttt{RuntimeException}（成为运行时异常），根据需求决定是否需要强制处理。②定义构造方法——通常提供多个构造方法，方便传递错误信息。至少应提供无参构造方法和带 String 参数的构造方法（用于传入错误描述信息），还可提供带 Throwable 参数的构造方法用于链式异常。③可选地，可以覆盖 \texttt{toString()}、\texttt{getMessage()}等方法提供更详细的异常信息。示例代码：
}
\vspace{-0.6cm}
\begin{lstlisting}[language=Java, basicstyle=\color{answerred}\ttfamily, keywordstyle=\color{blue}\bfseries, commentstyle=\color{green!50!black}, stringstyle=\color{red!70!black}, showstringspaces=false]
// 自定义受检查异常
public class InsufficientBalanceException extends Exception {
 public InsufficientBalanceException() {
 super();
 }
 public InsufficientBalanceException(String message) {
 super(message);
 }
 public InsufficientBalanceException(String message, Throwable cause) {
 super(message, cause);
 }
}
\end{lstlisting}
\vspace{-0.4cm}
\answercontinue{
使用自定义异常时，可以在需要的地方用 \texttt{throw}抛出，并在调用处用 try-catch 处理或继续向上抛出。}

\sq \texttt{String} 和 \texttt{StringBuilder}、\texttt{StringBuffer} 的区别是什么？
\answer{\texttt{String}、\texttt{StringBuilder}和 \texttt{StringBuffer}都是用于处理字符串的类，但它们在可变性和线程安全性上有重要区别：①String——不可变类，一旦创建，其内容不能被改变。每次对 String 进行修改（如拼接、替换）都会生成新的 String 对象，如果频繁操作会导致大量临时对象，影响性能。String 适用于内容不经常变化的场景（如常量、少量拼接）。②StringBuilder——可变类，可对其内容进行修改，不会产生新对象。提供了 append()、insert()、delete()等方法。StringBuilder 线程不安全（不同步），但性能最高。适用于单线程环境下大量字符串操作（如循环拼接）。③StringBuffer——可变类，与 StringBuilder 用法类似，但所有方法都用 \texttt{synchronized}修饰，是线程安全的。适用于多线程环境下操作字符串。性能比 StringBuilder 低。总结：操作少量字符串用 String；单线程大量操作用 StringBuilder；多线程大量操作用 StringBuffer。}

\sq 列举 Java 集合框架中常用的 List、Set、Map 接口的实现类，并简述它们的特点。
\answer{Java 集合框架（Collection Framework）提供了丰富的接口和实现类，常用的有三类：List、Set、Map。①List 接口（有序可重复）：实现类有 ArrayList（底层数组，查询快、增删慢，线程不安全）、LinkedList（底层双向链表，增删快、查询慢，线程不安全）、Vector（底层数组，线程安全，已过时）。②Set 接口（无序不可重复）：实现类有 HashSet（底层 HashMap，无序，允许 null，线程不安全）、LinkedHashSet（维护插入顺序，双向链表 + 哈希表）、TreeSet（底层 TreeMap，红黑树，元素自动排序，需实现 Comparable）。③Map 接口（键值对，键唯一）：实现类有 HashMap（底层数组 + 链表 + 红黑树，无序，允许 null 键值，线程不安全）、LinkedHashMap（维护插入顺序或访问顺序）、Hashtable（线程安全，不允许 null，已过时）、TreeMap（键自动排序，红黑树）。选择合适的集合类可提高程序效率和可维护性。}

\sq 什么是泛型？使用泛型有什么好处？
\answer{泛型（Generics）是 JDK 1.5 引入的重要特性，它允许在定义类、接口和方法时使用类型参数（Type Parameters），从而实现参数化类型。泛型的本质是将数据类型作为参数传递，使得同一套代码可以操作多种数据类型，同时保证类型安全。使用泛型的好处：①提高类型安全——在编译阶段就能检查类型错误，避免运行时出现 \texttt{ClassCastException}。例如，使用 \texttt{List<String>}明确指定列表只能存放 String 对象，插入其他类型会编译报错。②消除强制类型转换——从集合中取出元素时不需要显式强制转换，编译器自动插入正确的类型转换，代码更简洁。例如，\texttt{String s = list.get(0);}无需再写 \texttt{(String)}。③促进代码复用——泛型类、泛型方法可以适用于多种类型，减少重复代码。例如，\texttt{ArrayList<T>}可以存储任何类型。④增强可读性——泛型代码文档化效果更好，一看便知集合中存储的是什么类型。}

\sq 简述线程的生命周期，并说明线程创建的方式。
\answer{线程的生命周期是指线程从创建到消亡的整个过程，Java 中线程有六种状态（定义在 \texttt{Thread.State}枚举中）：①新建（New）——使用 \texttt{new Thread()}创建线程对象，尚未调用 \texttt{start()}。②就绪（Runnable）——调用 \texttt{start()}后，线程进入就绪队列等待 CPU 调度。③运行（Running）——线程获得 CPU 时间片，执行 \texttt{run()}方法中的代码。④阻塞（Blocked）——线程因等待获取锁而暂停执行。⑤等待（Waiting）——线程无限期等待，需被显式唤醒。⑥计时等待（Timed Waiting）——线程在指定时间内等待（如调用 \texttt{sleep()}、\texttt{wait(timeout)}）。⑦终止（Terminated）——线程的 \texttt{run()}方法执行完毕或抛出未捕获异常而结束。线程创建的方式：①继承 Thread 类——创建 Thread 子类，重写 \texttt{run()}方法，创建子类对象并调用 \texttt{start()}。②实现 Runnable 接口——实现 \texttt{run()}方法，创建 Thread 对象时将 Runnable 实例作为参数传入，调用 \texttt{start()}。③实现 Callable 接口——配合 FutureTask 使用，可返回结果和抛出异常。④使用线程池——通过 \texttt{ExecutorService}创建和管理线程。}

\sq 什么是线程同步？为什么要进行线程同步？synchronized 关键字可以怎么使用？
\answer{线程同步是指多个线程在访问共享资源时，通过某种机制协调它们的执行顺序，确保同一时刻只有一个线程能够访问临界资源，避免数据不一致问题。需要进行线程同步的原因：当多个线程并发访问同一共享变量（如银行账户余额、商品库存）时，如果不加控制，可能会出现"竞态条件"（Race Condition），导致数据错误。例如，两个线程同时取款，可能读取到相同余额，各自扣款后余额错误。synchronized 关键字是 Java 中最基本的线程同步机制，它能保证方法或代码块的原子性、可见性和有序性。synchronized 的三种使用方式：①修饰实例方法——同步当前实例对象（this），同一实例的多个线程不能同时执行该方法。②修饰静态方法——同步当前类的 Class 对象，对所有实例生效。③修饰代码块——指定锁对象，可以精确控制需要同步的代码范围，提高性能。例如：\texttt{synchronized (lockObject) { ... }}。除了 synchronized，还可使用 Lock 接口（如 ReentrantLock）、volatile 关键字、原子类等实现同步。}

\sq 简述 Java 中 I/O 流的分类，并举例说明字节流和字符流的区别。
\answer{Java I/O 流用于处理输入输出操作，根据不同的维度有多种分类方式：①按数据流向分为输入流（InputStream/Reader）和输出流（OutputStream/Writer）；②按处理单位分为字节流（以字节为单位，处理二进制数据）和字符流（以字符为单位，处理文本数据）；③按功能分为节点流（直接连接数据源）和处理流（包装其他流，提供额外功能，如缓冲、转换）。字节流和字符流的区别：①处理单位不同：字节流以 8 位字节为单位，字符流以 16 位 Unicode 字符为单位。②适用场景不同：字节流适合处理所有类型的数据（如图片、音频、视频、可执行文件等二进制文件）；字符流适合处理文本文件，可避免字符编码问题。③基类不同：字节流的抽象基类是 \texttt{InputStream}和 \texttt{OutputStream}；字符流的抽象基类是 \texttt{Reader}和 \texttt{Writer}。④典型实现：字节流如 \texttt{FileInputStream}、\texttt{FileOutputStream}；字符流如 \texttt{FileReader}、\texttt{FileWriter}。⑤缓冲区使用：字符流通常配合缓冲流（如 \texttt{BufferedReader}）提高效率。Java 提供了转换流（\texttt{InputStreamReader}/\texttt{OutputStreamWriter}）实现字节流与字符流之间的转换。}

\sq 什么是序列化？如何实现对象的序列化和反序列化？
\answer{序列化（Serialization）是指将对象的状态信息（成员变量数据）转换为可以存储或传输的字节序列的过程。反序列化（Deserialization）则是将字节序列恢复为对象的过程。序列化机制使得对象可以持久化到文件中，或通过网络在不同 JVM 之间传输。要实现序列化，需满足以下条件：①类必须实现 \texttt{Serializable}接口（这是一个标记接口，没有方法）。②类的所有属性都必须是可序列化的。如果某个属性不需要序列化，可以用 \texttt{transient}关键字修饰，该属性将被忽略。③可以定义 \texttt{serialVersionUID}（序列化版本号）来控制版本兼容性。序列化和反序列化的实现步骤：①创建 \texttt{ObjectOutputStream}对象，包装一个 \texttt{FileOutputStream}或其他输出流。②调用 \texttt{writeObject(对象)}方法将对象写入流。③创建 \texttt{ObjectInputStream}对象，包装对应的输入流。④调用 \texttt{readObject()}方法读取对象，返回值为 \texttt{Object}类型，需强制转换。示例代码：
}
\vspace{-0.6cm}
\begin{lstlisting}[language=Java, basicstyle=\color{answerred}\ttfamily, keywordstyle=\color{blue}\bfseries, commentstyle=\color{green!50!black}, stringstyle=\color{red!70!black}, showstringspaces=false]
// 序列化
ObjectOutputStream oos = new ObjectOutputStream(new FileOutputStream("obj.dat"));
oos.writeObject(student);
oos.close();
// 反序列化
ObjectInputStream ois = new ObjectInputStream(new FileInputStream("obj.dat"));
Student s = (Student) ois.readObject();
ois.close();
\end{lstlisting}
\vspace{-0.4cm}
\answercontinue{}

\sq 简述 JDBC 操作数据库的基本步骤。
\answer{JDBC（Java Database Connectivity）是 Java 程序连接和操作数据库的 API。使用 JDBC 操作数据库的基本步骤如下：①加载并注册数据库驱动——使用 \texttt{Class.forName("com.mysql.cj.jdbc.Driver")}动态加载驱动类（JDBC 4.0 后可以省略）。②建立数据库连接——通过 \texttt{DriverManager.getConnection(url, user, password)}获取 \texttt{Connection}对象，url 格式为 \texttt{jdbc:mysql://host:port/database}。③创建 Statement 对象——使用 \texttt{Connection}的 \texttt{createStatement()}创建 \texttt{Statement}对象，或使用 \texttt{prepareStatement()}创建预编译的 \texttt{PreparedStatement}对象（推荐，可防 SQL 注入）。④执行 SQL 语句——执行查询用 \texttt{executeQuery()}返回 \texttt{ResultSet}结果集；执行增删改用 \texttt{executeUpdate()}返回受影响行数。⑤处理结果集——对于查询操作，遍历 \texttt{ResultSet}获取数据。⑥关闭资源——依次关闭 \texttt{ResultSet}、\texttt{Statement}、\texttt{Connection}，释放数据库资源（最好在 finally 块中完成）。示例代码体现 try-with-resources 自动关闭资源的方式，可以简化代码。此外，实际开发中常使用连接池（如 HikariCP）提高性能。}

\begin{center}
 \textbf{—— 试卷结束，请仔细检查 ——}
\end{center}

\end{document}
